% DOC NO. RL-BRAND-001-PRINT · REV. A
%
% This file IS the document. Edit it directly - ripdoc compiles it
% as it stands and stamps the provenance line at build time.
%
%   ripdoc build --only docs-08-print
%
% Promoted from docs/08-print.md on 2026-09-07 by `ripdoc promote`.
\documentclass[fontpath=fonts/,logopath=logo/]{riposte-doc}
\usepackage{riposte-md}

\title{Print}
\author{Riposte Laboratories}
\date{}
\ripdocno{RL-BRAND-001-PRINT}
\riprev{A}
\ripstrapline{The brand was designed as printed matter, so it prints well — but screen RGB and press CMYK are different colour spaces and the accents are the ones that suffer.}
\riprunningtitle{Print}

\begin{document}

\ripmakehead

\riptoc


\ripsection{\ripmdhead{1 · The honest warning about CMYK}}

The CMYK values in \ripdocref{\ripmono{color-\allowbreak{}reference.\allowbreak{}md}}{RL-BRAND-001-COLOR-REFERENCE} and \ripmono{palettes/\allowbreak{}riposte.\allowbreak{}csv} are \textbf{naive device conversions}, computed arithmetically from RGB. They are a starting point to hand a printer. They are \textbf{not} a colour specification.

Three of our colours sit outside or near the edge of the CMYK gamut:

\begin{ripmdtable}{X[0.55,l]X[1.64,l]}
\ripmdth{COLOUR} & \ripmdth{PROBLEM ON PRESS} \\
\ripmono{teal} \ripmono{\#12B795} & Well outside CMYK gamut. Will print duller and darker than screen. The single biggest shift in the palette. \\
\ripmono{pink} \ripmono{\#F0477D} & Near the gamut edge. Loses fluorescence; drifts toward a flatter rose. \\
\ripmono{orange} \ripmono{\#FE9A0D} & Prints close, but can go muddy if the press pulls magenta. \\
\end{ripmdtable}

\textbf{Always ask for a physical proof} before a run of any size. Approve the proof, not the numbers.


\ripsection{\ripmdhead{2 · Spot colour (recommended for anything that matters)}}

For business cards, stickers, packaging, labels and signage, specify \textbf{spot inks}. The accents are the brand; a dull teal is a broken brand.

Starting points to confirm against a physical Pantone book under D50 light — \textbf{do not approve from a screen}:

\begin{ripmdtable}{X[0.55,l]X[0.57,l]X[2.01,l]}
\ripmdth{TOKEN} & \ripmdth{HEX} & \ripmdth{NEAREST PANTONE (COATED) — VERIFY PHYSICALLY} \\
\ripmono{ink} & \ripmono{\#1D1A17} & Black 6 C, or Neutral Black C \\
\ripmono{bone} & \ripmono{\#F6F1E7} & Warm Gray 1 C at low tint, or an uncoated natural stock \\
\ripmono{pink} & \ripmono{\#F0477D} & 191 C / 1915 C \\
\ripmono{orange} & \ripmono{\#FE9A0D} & 1375 C / 137 C \\
\ripmono{teal} & \ripmono{\#12B795} & 3395 C / 339 C \\
\end{ripmdtable}

\textbf{Bone should usually be the paper, not an ink.} A warm uncoated natural stock is closer to the brand than printing \ripmono{\#F6F1E7} onto white — and cheaper. See §5.


\ripsection{\ripmdhead{3 · Rule weights — the thing that breaks}}

The system is drawn in 2px screen rules. At print scale, translate as:

\begin{ripmdtable}{X[0.70,l]X[0.55,l]X[1.76,l]}
\ripmdth{SCREEN} & \ripmdth{PRINT} & \ripmdth{WHY} \\
\ripmono{1px} dashed divider & \textbf{0.5pt} & Below this it drops out \\
\ripmono{2px} solid structural rule & \textbf{1.5pt minimum}, 2pt preferred & \textbf{Do not go below 1pt.} The entire visual structure is these rules; if they vanish, the design vanishes \\
\ripmono{6px} accent top bar & \textbf{3mm} & {} \\
\ripmono{8px} accent left bar & \textbf{4mm} & {} \\
\end{ripmdtable}

Set all rules to \textbf{overprint off} and hold them at 100\% of a single ink where possible — a 1.5pt rule built from four-colour black will show registration fringing.


\ripsection{\ripmdhead{4 · Type on paper}}

\begin{ripmdtable}{X[0.55,l]X[0.73,l]X[1.82,l]}
\ripmdth{ROLE} & \ripmdth{SIZE} & \ripmdth{NOTES} \\
Body & \textbf{9.5–10pt} / 15pt leading & JetBrains Mono is wide; don't go below 9pt \\
\ripmono{md} secondary & 8.5–9pt & {} \\
Labels, chrome & 7–7.5pt, \ripmono{.1em} tracked & The tracking matters more in print, not less \\
h3 & 14–16pt & {} \\
h2 & 18–22pt & {} \\
h1 / page head & 28–36pt & {} \\
Display & 60pt+ & {} \\
\end{ripmdtable}

\begin{ripbullets}
\item \textbf{Embed the font.} JetBrains Mono is OFL — embedding and redistribution are permitted.
\item Weight 400 for body, 700 for headings, 800 for display and stat figures.
\item Ink text should be \textbf{100\% K plus a touch of warmth} (e.g. \ripmono{C0 M10 Y21 K89} from the conversion table) — not registration black, which cracks on folds and shows through.
\item Reversed type (bone on ink) needs a \textbf{1pt minimum stroke weight}; at 7pt, tracked uppercase, thin strokes fill in with ink spread. Bump to 8pt reversed.
\end{ripbullets}


\ripsection{\ripmdhead{5 · Paper}}

The whole visual language is "printed on bone paper". Choose stock accordingly:

\begin{ripbullets}
\item \textbf{Uncoated, warm white or natural.} Not bright white, not coated gloss.
\item Target something close to \ripmono{\#F6F1E7} in the stock itself. Then \ripmono{bone} costs no ink and looks right by definition.
\item Weights: 120–170 gsm text, 300–350 gsm cards and covers.
\item \textbf{Avoid gloss.} The brand has no shine — no blur, no glow, no gradients. A soft-touch or uncoated finish is on-message; gloss lamination is not.
\item Recycled and post-consumer stock is on-message in the most literal possible way. Say so in the colophon when you use it.
\end{ripbullets}


\ripsection{\ripmdhead{6 · Layout}}

\begin{ripmdgrid}{X[0.55,l]X[1.78,l]}
Margins & \textbf{15mm minimum, 18mm preferred} \\
Bleed & 3mm, with all bands and full-bleed fields extended into it \\
Spacing & \ripmono{--space-4} → 4mm · \ripmono{--space-6} → 8mm · \ripmono{--space-8} → 15mm \\
Radius & \textbf{0.} Do not let a template round a corner \\
Bands & Checker 26px → \textbf{7mm} · Harlequin 38px → \textbf{10mm} \\
\end{ripmdgrid}

Keep the \ripmono{.pagehead} tri-band and the colophon footer. A printed Riposte document should be recognisable as the same object as the web page.


\ripsection{\ripmdhead{7 · Web-to-print (\NoCaseChange{\ripmono{@media print}})}}

\ripmono{brand/\allowbreak{}riposte-\allowbreak{}brand.\allowbreak{}css} ships a print block that:

\begin{ripbullets}
\item drops to true black on white (ink-on-bone at 15.39:1 is lovely on screen; on a home laser printer, bone becomes a grey wash)
\item hides \ripmono{.marquee}, \ripmono{.checker} and \ripmono{.harlequin} — decoration that wastes toner
\item underlines links
\end{ripbullets}

If you're generating a PDF for actual press (not a desktop printer), \textbf{override that block} and keep the real colours and bands. It exists for "someone hit Ctrl-P", not for production.


\ripsection{\ripmdhead{8 · Handing off to a printer}}

Send:

\begin{ripnumbers}
\item The artwork as \textbf{PDF/X-4}, fonts embedded, 3mm bleed, crop marks
\item \ripmdpath{\ripmono{palettes/\allowbreak{}riposte.\allowbreak{}ase}} — loads in Illustrator, InDesign, Affinity
\item \ripmdpath{\ripmono{palettes/\allowbreak{}riposte.\allowbreak{}csv}} — every value in every notation
\item This page
\end{ripnumbers}

Ask for:

\begin{ripbullets}
\item A \textbf{physical proof} on the actual stock
\item \textbf{Spot inks} for the three accents where budget allows
\item Confirmation that \textbf{rules are holding at 1.5pt} and are not being auto-thinned
\end{ripbullets}


\ripsection{\ripmdhead{9 · Checklist}}

\begin{riptodolist}
\riptodo Physical proof requested and approved on the real stock
\riptodo Structural rules ≥ 1.5pt
\riptodo Fonts embedded (OFL permits it)
\riptodo Radius 0 everywhere
\riptodo Bleed 3mm, bands extended into it
\riptodo Ink text is warm black, not registration black
\riptodo Reversed type ≥ 8pt
\riptodo Colophon present
\riptodo Teal checked on the proof specifically — it's the one that shifts most
\end{riptodolist}

\ripend

\end{document}
